\paragraph{5. Error-Prone Average-Case Easiness of NP.}
In this section, we show the error-prone average-case easiness of $\mathsf{NP}$.
\paragraph{Theorem 15.}
With probability 1 over the choice of $\mathcal O_a$, $\mathsf{DistNP}^{\mathcal O_a}\subseteq\mathsf{HeurP}^{\mathcal O_a}$ holds.
First, we introduce several notations. For each choice of $\mathcal O$, we define the oracle $\mathcal A^*$ in the same manner as $\mathcal A$ except we apply $\rho_{1,\leq i_T-1}\rho^*_{1,i_T},\ldots,\rho_{T,\leq i_T-1}\rho^*_{T,i_T}$ to $\phi$ instead of $\rho_{1,\leq i_T},\ldots,\rho_{T,\leq i_T}$. Note that $\mathcal A^*$ executes a given nondeterministic machine $M$ with access to $\mathcal A$ (rather than $\mathcal A^*$) to construct the corresponding DNF $\phi$. We can verify that $\mathcal A^*$ is well-defined (i.e., not circular) in the same manner as Proposition 14.
Now, we show that $\mathcal A$ and $\mathcal A^*$ do not differ considerably.
\paragraph{Lemma 16.}
For each input $(M,x,1^{T^{2^c}})$ to $\mathcal A$, we have that
\[\Pr_{\mathcal O}[\mathcal A(M,x,1^{T^{2^c}})\neq\mathcal A^*(M,x,1^{T^{2^c}})]=O(T^{-4}).\]
\begin{proof}
Let $i:=i_T=(1/c)\log\log T$. For all $n\leq T$ and $j\leq i-1$, we fix $\rho^*_{n,j}$, $S_{n,j}$, and $\rho^*_{n,i}$ arbitrarily; let $C_T$ denote this condition. Notice that the DNF formula $\phi_{C_T}$ constructed by $\mathcal A$ and $\mathcal A^*$ is determined only by $C_T$ because all answers to recursive calls for $\mathcal A$ are determined by $C_T$ as in Proposition 14. Let $\phi'_{C_T}=\phi_{C_T}|_{\rho_{1,\leq i-1},\ldots,\rho_{T,\leq i-1}}$. Then, the value of $\mathcal A(M,x,1^{T^{2^c}})$ (resp. $\mathcal A^*(M,x,1^{T^{2^c}})$) is determined by $\phi'_{C_T}|_{\rho_{1,i},\ldots,\rho_{T,i}}$ (resp. $\phi'_{C_T}|_{\rho^*_{1,i},\ldots,\rho^*_{T,i}}$).
For any DNF formula $\phi$ and a restriction $\rho$, there are the following three cases: (i) $\phi|_\rho\equiv0$, (ii) $\phi|_\rho\equiv1$, or (iii) $\phi|_\rho$ does not become a constant (we write this case as $\phi|_\rho\equiv*$). Following this case analysis, there exist $3^2=9$ cases on $(\phi'_{C_T}|_{\rho^*_{1,i},\ldots,\rho^*_{T,i}},\phi'_{C_T}|_{\rho_{1,i},\ldots,\rho_{T,i}})$. However, since each $\rho_{n,i}$ is a subrestriction of $\rho^*_{n,i}$ (i.e., $\rho_{n,i}$ assigns values only to variables that are also assigned by $\rho^*_{n,i}$), the following 4 cases do not occur: $(0,1)$, $(1,0)$, $(*,0)$, and $(*,1)$. Further, in the cases of $(0,0)$, $(1,1)$, $(*,*)$, and $(*,0)$, the answers by $\mathcal A^*$ and $\mathcal A$ do not differ because $\mathcal A^*$ and $\mathcal A$ return 0 in the case of $*$. Thus, we only need to consider the case of $(1,*)$.
By the aforementioned argument, the probability in the lemma is expressed as follows:
\begin{align*}
&\Pr_{\mathcal O}[\mathcal A(M,x,1^{T^{2^c}})\neq\mathcal A^*(M,x,1^{T^{2^c}})]\\
&=\mathbb E_{C_T}\!\left[\Pr_{S_{1,i},\ldots,S_{T,i}}\!\left[\phi'_{C_T}|_{\rho_{1,i},\ldots,\rho_{T,i}}\equiv*\ \middle|\ C_T,\phi'_{C_T}|_{\rho^*_{1,i},\ldots,\rho^*_{T,i}}\equiv1\right]\right].
\end{align*}
To bound the probability in the right-hand side for each condition, we consider the case in which $\phi'_{C_T}|_{\rho^*_{1,i},\ldots,\rho^*_{T,i}}\equiv1$. Then, we can select a term $\tau$ in $\phi'_{C_T}$ such that $\tau|_{\rho^*_{1,i},\ldots,\rho^*_{T,i}}\equiv1$. For each $n\leq T$, define $Z_n\subseteq\{0,1\}^n$ as
\[Z_n=\{z\in\{0,1\}^n:\exists(x,\ell)\in\{0,1\}^n\times[n]\text{ s.t. a variable }F_{z,x,\ell}\text{ is contained in }\tau\}.\]
Since $\phi'_{C_T}$ is a $T$-DNF, $|Z_n|\leq T$ for each $n\leq T$. Furthermore, for any $n\in\mathbb N$ such that $i_{\max}(n)\leq i-1$, we have that $Z_n=\varnothing$ because all such variables must be assigned by $\rho_{n,\leq i-1}$. If $\phi'_{C_T}|_{\rho_{1,i},\ldots,\rho_{T,i}}\equiv*$, then $\tau|_{\rho_{1,i},\ldots,\rho_{T,i}}\equiv*$ must hold. This event occurs only if $\bigcup_{n\leq T}(S_{n,i}\cap Z_n)\neq\varnothing$ holds. Since each $S_{n,i}$ is selected from $S_{n,i-1}$ uniformly at random, this probability is bounded above by
\begin{align*}
\Pr\left[\bigcup_{n\leq T}(S_{n,i}\cap Z_n)\neq\varnothing\right]
&\leq\sum_{n\leq T}\Pr[S_{n,i}\cap Z_n\neq\varnothing]\\
&\leq\sum_{n\leq T:i_{\max}(n)\geq i}|Z_n|\cdot|S_{n,i}|/|S_{n,i-1}|\\
&\leq O(T)\cdot\sum_{n:t^{-1}(T)\leq n\leq T}p(n)\\
&\leq O(T^2\cdot t(t^{-1}(T))^{-6})=O(T^{-4}).
\end{align*}
Thus, we conclude that
\begin{align*}
&\Pr_{\mathcal O}[\mathcal A(M,x,1^{T^{2^c}})\neq\mathcal A^*(M,x,1^{T^{2^c}})]\\
&=\mathbb E_{C_T}\!\left[\Pr_{S_{1,i},\ldots,S_{T,i}}\!\left[\phi'_{C_T}|_{\rho_{1,i},\ldots,\rho_{T,i}}\equiv*\ \middle|\ C_T,\phi'_{C_T}|_{\rho^*_{1,i},\ldots,\rho^*_{T,i}}\equiv1\right]\right]\leq O(T^{-4}).
\end{align*}
\end{proof}
Furthermore, we can show the average-case easiness of $\mathsf{NP}$ under the oracle access to $\mathcal A^*$ (instead of $\mathcal A$). This part is similar to in [25,20]. Thus, we defer the proof to Appendix A.
\paragraph{Lemma 17 ([25,20]).}
Let $M$ be a $t_M(n)$-time nondeterministic oracle machine and $S$ be a randomized polynomial-time oracle sampling machine. We assume that $S(1^n)$ takes at most $t_S(n)$ time to generate an instance of length $n$. Then, the following event occurs with probability 1 over the choice of $\mathcal O$: for any $n\in\mathbb N$,
\[\Pr_{x\leftarrow S^{\mathcal O}(1^n)}[M^{\mathcal O}(x)\neq\mathcal A^*(M,x,1^{T^{2^c}})]\leq O(n^{-4}),\]
where $T=\max\{n^{2^c},t_M(n)^{2^c},t_S(n)^{2^c}\}$.
Now, we derive the average-case easiness of $\mathsf{NP}$ from Lemmas 16 and 17.
\begin{proof}[Proof of Theorem 15]
We consider an arbitrary distributional $\mathsf{NP}$ problem $(L,D)$ and assume that $L$ is specified by a $t_M(n)$-time nondeterministic oracle machine $M$, and $D$ is specified by a randomized $t_S(n)$-time oracle sampling machine $S$. Let $T=\max\{n^{2^c},t_M(n)^{2^c},t_S(n)^{2^c}\}$ as in Lemma 17.
We construct an error-prone heuristic algorithm $B$ for $(L,D)$ as follows: On input $x\in\{0,1\}^n$, $B$ queries $(M,x,1^{T^{2^c}})$ to $\mathcal A$ and returns the same answer. In the following, we will verify that the error probability of $B$ (over the choice of $\mathcal O$ and $S^{\mathcal O}(1^n)$) is bounded above by $O(n^{-4})$ for each input size $n$. Then, by applying Markov's inequality and the Borel--Cantelli lemma, the error probability of $B$ is bounded above by $n^{-2}$ for all sufficiently large $n$ with probability 1 over the choice of $\mathcal O$. Since the number of tuples $(M,S)$ is countable, we can conclude that all distributional $\mathsf{NP}$ problems have an error-prone heuristic algorithm with error probability at most $n^{-2}$ with probability 1 over the choice of $\mathcal O$. Based on the argument in [24, Proposition 3], this is sufficient for the statement that $\mathsf{DistNP}^{\mathcal O}\subseteq\mathsf{HeurP}^{\mathcal O}$.
Therefore, it is sufficient to show that, for any $n\in\mathbb N$,
\[\Pr_{\mathcal O,x\leftarrow S^{\mathcal O}(1^n)}[M^{\mathcal O}(x)\neq\mathcal A(M,x,1^{T^{2^c}})]\leq O(n^{-4}).\]
Obviously, this event occurs only if (i) $\mathcal A(M,x,1^{T^{2^c}})\neq\mathcal A^*(M,x,1^{T^{2^c}})$ or (ii) $M^{\mathcal O}(x)\neq\mathcal A^*(M,x,1^{T^{2^c}})$ occur. By Lemmas 16 and 17 and the union bound, we have
\begin{align*}
&\Pr_{\mathcal O,x\leftarrow S^{\mathcal O}(1^n)}[M^{\mathcal O}(x)\neq\mathcal A(M,x,1^{T^{2^c}})]\\
&\leq\Pr_{\mathcal O,x\leftarrow S^{\mathcal O}(1^n)}[\mathcal A(M,x,1^{T^{2^c}})\neq\mathcal A^*(M,x,1^{T^{2^c}})]\\
&\quad+\Pr_{\mathcal O,x\leftarrow S^{\mathcal O}(1^n)}[M^{\mathcal O}(x)\neq\mathcal A^*(M,x,1^{T^{2^c}})]\\
&=O(T^{-4})+O(n^{-4})=O(n^{-4})+O(n^{-4})=O(n^{-4}).
\end{align*}
\end{proof}
